Airtel Uganda Limited relies on Average Revenue Per User (ARPU)-banded segmentation to deliver voice bundle promotions, yielding conversion rates below 3 percent and wasting an estimated 20 percent of the marketing budget on untargeted communications. The absence of a personalised, data-driven recommendation system represents both a commercial inefficiency and an academic gap this study addresses.

This study developed and validated a machine learning recommendation framework using real transactional data from 1,458,900 Airtel Uganda prepaid subscribers (February--April 2026). Four classification models --- Logistic Regression, Decision Tree, Random Forest, and XGBoost --- were evaluated following a preprocessing pipeline of median imputation, capped SMOTE oversampling, and behavioural feature engineering. Three statistical tests were applied to evaluate relationships between customer attributes and bundle segment: Spearman Rank Correlation to measure the strength and direction of monotonic associations between numeric features and the target; one-way ANOVA to test whether feature means differ significantly across the eight value segment groups; and Chi-Square to assess the independence of categorical features from the target variable.

Statistical tests confirmed 34 of 35 features significantly associated with bundle segment ($p < 0.05$); UPGRADE\_RATE was the only non-significant feature. Feature importance analysis identified TOTAL\_REVENUE (18.01\%), SPEND\_PER\_MINUTE (6.57\%), CALLS (5.97\%), VOICE\_INTENSITY (5.87\%), and MOBILE\_MONEY (4.64\%) as the top five predictors. XGBoost achieved the highest F1-Score of 0.661 on the 8-class hold-out test set --- a 5.4-fold improvement over a random baseline.

A simulated pilot projected a four- to five-fold improvement in conversion rate (3\% to 12--15\%) and a 9 percent ARPU uplift. Three strategic recommendations were formulated: a VOICE\_INTENSITY-triggered recommendation engine, enriched segmentation incorporating device type and purchase channel, and a formal A/B test with model governance. This is the first empirically validated ML bundle recommendation framework built on real Ugandan telecom data.

\textbf{Keywords:} Machine Learning, Telecommunications, Voice bundle recommendation, Customer behavioural attributes, Predictive Analytics.
